% Part: first-order-logic
% Chapter: tableaux
% Section: proving-things-quant

\documentclass[../../../include/open-logic-section]{subfiles}

\begin{document}

\olfileid{fol}{tab}{prq}

\olsection{పరిమాణీకరణికాలతో \usetoken{P}{tableau}}

\begin{ex}
పరిమాణీకరణికాలతో వ్యవహరించేటప్పుడు ఐగెన్ చరరాశి షరతును
ఉల్లంఘించకుండా చూసుకోవాలి; దానికోసం కొన్ని నిగమనాలను చేసే క్రమాన్ని
కొన్నిసార్లు మార్చి ప్రయత్నించాలి. సాధారణంగా ఐగెన్ చరరాశి షరతుకు
లోబడే నియమాలను ముందుగా చూసుకోవడం సహాయపడుతుంది (పూర్తయిన
!!{tableau}లో అవి పై భాగంలో ఉంటాయి).

!!{sentence} $\lexists[x][\lnot !A(x)] \lif \lnot
\lforall[x][!A(x)]$ కోసం ఒక !!a{tableau}ను ఎలా ఇస్తామో చూద్దాం.
ఎప్పటిలాగే పరికల్పనను నమోదు చేస్తూ మొదలుపెడతాం:
\begin{oltableau}
[\sFmla{\False}{\lexists[x][\lnot \formula{A}(x)]\lif \lnot
    \lforall[x][\formula{A}(x)]}, just=\TAss]
\end{oltableau}
!!{main operator} $\lif$ కనుక $\TRule{\False}{\lif}$ను
వర్తింపజేస్తాం:
\begin{oltableau}
[\sFmla{\False}{\lexists[x][\lnot \formula{A}(x)]\lif \lnot
    \lforall[x][\formula{A}(x)]}, just=\TAss, checked
  [\sFmla{\True}{\lexists[x][\lnot \formula{A}(x)]},
    just={\TRule{\False}{\lif}[1]}
    [\sFmla{\False}{\lnot\lforall[x][\formula{A}(x)]},
      just={\TRule{\False}{\lif}[1]}]
  ]
]
\end{oltableau}
తర్వాత చూసుకోవాల్సింది వరుస~$2$. $\TRule{\True}{\lexists}$ను
ఉపయోగిస్తాం. దీనికి కొత్త !!{constant} కావాలి; ఇప్పటివరకు ఏ
!!{constant}s కనిపించలేదు కనుక ఏదైనా ఒకదాన్ని, ఉదాహరణకు~$a$ను,
ఎంచుకోవచ్చు.
\begin{oltableau}
[\sFmla{\False}{\lexists[x][\lnot \formula{A}(x)]\lif \lnot
    \lforall[x][\formula{A}(x)]}, just=\TAss, checked
  [\sFmla{\True}{\lexists[x][\lnot \formula{A}(x)]},
    just={\TRule{\False}{\lif}[1]}, checked
    [\sFmla{\False}{\lnot\lforall[x][\formula{A}(x)]},
      just={\TRule{\False}{\lif}[1]}
      [\sFmla{\True}{\lnot\formula{A}(a)}, just={\TRule{\True}{\lexists}[2]}]
    ]
  ]
]
\end{oltableau}
ఇప్పుడు వరుస~$3$కు $\TRule{\False}{\lnot}$ను వర్తింపజేస్తాం:
\begin{oltableau}
[\sFmla{\False}{\lexists[x][\lnot \formula{A}(x)]\lif \lnot
    \lforall[x][\formula{A}(x)]}, just=\TAss, checked
  [\sFmla{\True}{\lexists[x][\lnot \formula{A}(x)]},
    just={\TRule{\False}{\lif}[1]}, checked
    [\sFmla{\False}{\lnot\lforall[x][\formula{A}(x)]},
      just={\TRule{\False}{\lif}[1]}, checked
      [\sFmla{\True}{\lnot\formula{A}(a)}, just={\TRule{\True}{\lexists}[2]}
        [\sFmla{\True}{\lforall[x][\formula{A}(x)]},
          just = {\TRule{\False}{\lnot}[3]}
        ]
      ]
    ]
  ]
]
\end{oltableau}
వరుస~$4$కు $\TRule{\True}{\lnot}$ను, ఆ తర్వాత వరుస~$5$కు
$\TRule{\True}{\lforall}$ను వర్తింపజేసి సంవృత !!{tableau}ను
పొందుతాం.
\begin{oltableau}
[\sFmla{\False}{\lexists[x][\lnot \formula{A}(x)]\lif \lnot
    \lforall[x][\formula{A}(x)]}, just=\TAss, checked
  [\sFmla{\True}{\lexists[x][\lnot \formula{A}(x)]},
    just={\TRule{\False}{\lif}[1]}, checked
    [\sFmla{\False}{\lnot\lforall[x][\formula{A}(x)]},
      just={\TRule{\False}{\lif}[1]}, checked
      [\sFmla{\True}{\lnot\formula{A}(a)}, just={\TRule{\True}{\lexists}[2]}
        [\sFmla{\True}{\lforall[x][\formula{A}(x)]},
          just = {\TRule{\False}{\lnot}[3]}
          [\sFmla{\False}{\formula{A}(a)}, just={\TRule{\True}{\lnot}[4]}
            [\sFmla{\True}{\formula{A}(a)},
              just={\TRule{\True}{\lforall}[5]}, close
            ]
          ]
        ]
      ]
    ]
  ]
]
\end{oltableau}
\end{ex}

\begin{ex}
కింది సమితికి ఒక !!a{tableau}ను ఎలా ఇస్తామో చూద్దాం:
\[
\sFmla{\False}{\lexists[x][!C(x,b)]},
\sFmla{\True}{\lexists[x][(!A(x) \land !B(x))]},
\sFmla{\True}{\lforall[x][(!B(x) \lif !C(x,b))]}.
\]
ఎప్పటిలాగే పరికల్పనలతో మొదలుపెడతాం:
\begin{oltableau}
  [\sFmla{\False}{\lexists[x][\formula{C}(x,b)]}, just=\TAss
    [\sFmla{\True}{\lexists[x][(\formula{A}(x) \land \formula{B}(x))]},
      just=\TAss
      [\sFmla{\True}{\lforall[x][(\formula{B}(x) \lif \formula{C}(x,b))]},
        just=\TAss]
    ]
  ]
\end{oltableau}
ఐగెన్ చరరాశి షరతు గల నియమాన్ని ఎప్పుడూ ముందుగా
వర్తింపజేయాలి; ఇక్కడ వరుస~$2$కు $\TRule{\True}{\lexists}$ను
వర్తింపజేయాలి. పరికల్పనల్లో !!{constant}~$b$ ఉంది కనుక మరొకదాన్ని
వాడాలి; మళ్లీ~$a$నే ఎంచుకుందాం.
\begin{oltableau}
  [\sFmla{\False}{\lexists[x][\formula{C}(x,b)]}, just=\TAss
    [\sFmla{\True}{\lexists[x][(\formula{A}(x) \land \formula{B}(x))]},
      just=\TAss, checked
      [\sFmla{\True}{\lforall[x][(\formula{B}(x) \lif \formula{C}(x,b))]},
        just=\TAss
        [\sFmla{\True}{\formula{A}(a) \land \formula{B}(a)},
          just={\TRule{\True}{\lexists}[2]}]
      ]
    ]
  ]
\end{oltableau}
ఇప్పుడు వరుస~$1$కు $\TRule{\False}{\lexists}$ను లేదా వరుస~$3$కు
$\TRule{\True}{\lforall}$ను వర్తింపజేస్తే, $x$ స్థానంలో ఏ పదం~$t$ను
ప్రతిస్థాపించాలో నిర్ణయించాలి. ఈ నియమాలకు ఐగెన్ చరరాశి షరతు లేదు
కనుక మనకు నచ్చిన పదాన్ని ఎంచుకోవచ్చు. కొన్ని సందర్భాల్లో వేర్వేరు
$t$లతో నియమాన్ని అనేకసార్లు వర్తింపజేయాల్సి రావచ్చు. అయితే సాధారణ
నియమంగా, !!{tableau}లో ఇప్పటికే కనిపించే పదాల్లో ఒకదాన్ని---ఇక్కడ
$a$ లేదా~$b$ను---ఎంచుకోవడం ప్రయోజనకరం; ఈ సందర్భంలో $a$ సంవృత శాఖను
ఇచ్చే అవకాశం ఎక్కువని ఊహించవచ్చు.
\begin{oltableau}
  [\sFmla{\False}{\lexists[x][\formula{C}(x,b)]}, just=\TAss
    [\sFmla{\True}{\lexists[x][(\formula{A}(x) \land \formula{B}(x))]},
      just=\TAss, checked
      [\sFmla{\True}{\lforall[x][(\formula{B}(x) \lif \formula{C}(x,b))]}, just=\TAss
        [\sFmla{\True}{\formula{A}(a) \land \formula{B}(a)}, just={\TRule{\True}{\lexists}[2]}
          [\sFmla{\False}{\formula{C}(a,b)}, just={\TRule{\False}{\lexists}[1]}
            [\sFmla{\True}{\formula{B}(a) \lif \formula{C}(a,b)},
              just={\TRule{\True}{\lforall}[3]}
            ]
          ]
        ]
      ]
    ]
  ]
\end{oltableau}
వరుసలు $1$, $3$లోని !!{signed formula}sకు తనిఖీ గుర్తు పెట్టం;
వాటిని మళ్లీ ఉపయోగించాల్సి రావచ్చు. ఇప్పుడు వరుస~$4$కు
$\TRule{\True}{\land}$ను వర్తింపజేద్దాం:
\begin{oltableau}
  [\sFmla{\False}{\lexists[x][\formula{C}(x,b)]}, just=\TAss
    [\sFmla{\True}{\lexists[x][(\formula{A}(x) \land \formula{B}(x))]},
      just=\TAss, checked
      [\sFmla{\True}{\lforall[x][(\formula{B}(x) \lif \formula{C}(x,b))]}, just=\TAss
        [\sFmla{\True}{\formula{A}(a) \land \formula{B}(a)},
          just={\TRule{\True}{\lexists}[2]}, checked
          [\sFmla{\False}{\formula{C}(a,b)}, just={\TRule{\False}{\lexists}[1]}
            [\sFmla{\True}{\formula{B}(a) \lif \formula{C}(a,b)},
              just={\TRule{\True}{\lforall}[3]}
              [\sFmla{\True}{\formula{A}(a)}, just={\TRule{\True}{\land}[4]}
                [\sFmla{\True}{\formula{B}(a)}, just={\TRule{\True}{\land}[4]}
                ]
              ]
            ]
          ]
        ]
      ]
    ]
  ]
\end{oltableau}
ఇప్పుడు వరుస~$6$కు $\TRule{\True}{\lif}$ను వర్తింపజేస్తే
!!{tableau} సంవృతమవుతుంది:
\begin{oltableau}
  [\sFmla{\False}{\lexists[x][\formula{C}(x,b)]}, just=\TAss
    [\sFmla{\True}{\lexists[x][(\formula{A}(x) \land \formula{B}(x))]},
      just=\TAss, checked
      [\sFmla{\True}{\lforall[x][(\formula{B}(x) \lif \formula{C}(x,b))]},
        just=\TAss
        [\sFmla{\True}{\formula{A}(a) \land \formula{B}(a)},
          just={\TRule{\True}{\lexists}[2]}, checked
          [\sFmla{\False}{\formula{C}(a,b)},
            just={\TRule{\False}{\lexists}[1]}
            [\sFmla{\True}{\formula{B}(a) \lif \formula{C}(a,b)},
              just={\TRule{\True}{\lforall}[3]}, checked
              [\sFmla{\True}{\formula{A}(a)},
                just={\TRule{\True}{\land}[4]}
                [\sFmla{\True}{\formula{B}(a)},
                  just={\TRule{\True}{\land}[4]}
                  [\sFmla{\False}{\formula{B}(a)},
                    just={\TRule{\True}{\lif}[6]},close
                  ]
                  [\sFmla{\True}{\formula{C}(a,b)},
                    just={\TRule{\True}{\lif}[6]},close
                  ]
                ]
              ]
            ]
          ]
        ]
      ]
    ]
  ]
\end{oltableau}


\end{ex}

\begin{ex}
కింది సమితికి ఒక !!a{tableau}ను నిర్మిస్తాం:
\[
\sFmla{\True}{\lforall[x][!A(x)]}, \sFmla{\True}{\lforall[x][!A(x)]
  \lif \lexists[y][!B(y)]}, \sFmla{\True}{\lnot\lexists[y][!B(y)]}.
\]
ఎప్పటిలాగే పరికల్పనలను రాస్తాం:
\begin{oltableau}
  [\sFmla{\True}{\lforall[x][\formula{A}(x)]}, just=\TAss
    [\sFmla{\True}{\lforall[x][\formula{A}(x)] \lif
        \lexists[y][\formula{B}(y)]}, just=\TAss
      [\sFmla{\True}{\lnot\lexists[y][\formula{B}(y)]}, just=\TAss
      ]
    ]
  ]
\end{oltableau}
మొదట వరుస~$3$కు $\TRule{\True}{\lnot}$ నియమాన్ని వర్తింపజేస్తాం.
“ఐగెన్ చరరాశి షరతులు గల నియమాలను ఎప్పుడూ ముందుగా వర్తింపజేయాలి”
అనే నియమానికి పర్యవసానం: “ఐగెన్ చరరాశి షరతులు లేని పరిమాణీకరణిక
నియమాలు అవసరమయ్యే వరకు వాయిదా వేయాలి.” శాఖగా విడగొట్టే నియమాలనూ
వాయిదా వేయండి.
\begin{oltableau}
  [\sFmla{\True}{\lforall[x][\formula{A}(x)]}, just=\TAss
    [\sFmla{\True}{\lforall[x][\formula{A}(x)] \lif
        \lexists[y][\formula{B}(y)]}, just=\TAss
      [\sFmla{\True}{\lnot\lexists[y][\formula{B}(y)]}, just=\TAss, checked
        [\sFmla{\False}{\lexists[y][\formula{B}(y)]}, just={\TRule{\True}{\lnot}[3]}]
      ]
    ]
  ]
\end{oltableau}
కొత్త వరుస~$4$కు ఐగెన్ చరరాశి షరతు లేని పరిమాణీకరణిక నియమం
$\TRule{\False}{\lexists}$ అవసరం. కనుక దాన్ని వాయిదా వేసి,
వరుస~$2$కు $\TRule{\True}{\lif}$ను ఉపయోగిస్తాం.
\begin{oltableau}
  [\sFmla{\True}{\lforall[x][\formula{A}(x)]}, just=\TAss
    [\sFmla{\True}{\lforall[x][\formula{A}(x)] \lif
        \lexists[y][\formula{B}(y)]}, just=\TAss, checked
      [\sFmla{\True}{\lnot\lexists[y][\formula{B}(y)]}, just=\TAss, checked
        [\sFmla{\False}{\lexists[y][\formula{B}(y)]},
          just={\TRule{\True}{\lnot}[3]},
          [\sFmla{\False}{\lforall[x][\formula{A}(x)]}, just={\TRule{\True}{\lif}[2]}]
          [\sFmla{\True}{\lexists[y][\formula{B}(y)]}, just={\TRule{\True}{\lif}[2]}]
        ]
      ]
    ]
  ]
\end{oltableau}
రెండు కొత్త !!{signed formula}sకూ ఐగెన్ చరరాశి షరతులు గల నియమాలు
అవసరం; కాబట్టి తర్వాత వాటినే వర్తింపజేయాలి:
\begin{oltableau}
  [\sFmla{\True}{\lforall[x][\formula{A}(x)]}, just=\TAss
    [\sFmla{\True}{\lforall[x][\formula{A}(x)] \lif
        \lexists[y][\formula{B}(y)]}, just=\TAss, checked
      [\sFmla{\True}{\lnot\lexists[y][\formula{B}(y)]}, just=\TAss, checked
        [\sFmla{\False}{\lexists[y][\formula{B}(y)]},
          just={\TRule{\True}{\lnot}[3]}
          [\sFmla{\False}{\lforall[x][\formula{A}(x)]}, just={\TRule{\True}{\lif}[2]},checked
            [\sFmla{\False}{\formula{A}(b)}, just={\TRule{\False}{\lforall}[5]}]
          ]
          [\sFmla{\True}{\lexists[y][\formula{B}(y)]}, just={\TRule{\True}{\lif}[2]},checked
            [\sFmla{\True}{\formula{B}(c)}, just={\TRule{\True}{\lexists}[5]}]
          ]
        ]
      ]
    ]
  ]
\end{oltableau}
శాఖలను సంవృతం చేయడానికి వరుసలు $1$,~$4$లోని !!{signed formula}sను
ఉపయోగించాలి. వాటికి సంబంధించిన నియమాలకు
(\TRule{\True}{\lforall}, \TRule{\False}{\lexists}) ఐగెన్ చరరాశి
షరతులు లేవు; కనుక తగిన పదాలను స్వేచ్ఛగా ఎంచుకోవచ్చు. ఇక్కడ అవి
వరుసగా $b$,~$c$.
\sourcecorrection{OLTETAB-003}{ఆంగ్ల మూలం ఈ చివరి దశకు అవసరమైన
చిహ్నిత సూత్రాలను “వరుసలు $1$, $3$”లో ఉన్నవని చెప్పింది. కానీ
వరుస~$3$లోని నిషేధానికి ఇప్పటికే తనిఖీ గుర్తు ఉంది; అవసరమైన
$\sFmla{\False}{\lexists[y][!B(y)]}$ వరుస~$4$లో ఉందని ముందరి,
తదుపరి టాబ్లోలు రెండూ చూపుతున్నాయి. అందువల్ల వరుస సూచనను $1$,~$4$గా
సరిచేశాం.}
\begin{oltableau}
  [\sFmla{\True}{\lforall[x][\formula{A}(x)]}, just=\TAss
    [\sFmla{\True}{\lforall[x][\formula{A}(x)] \lif
        \lexists[y][\formula{B}(y)]}, just=\TAss, checked
      [\sFmla{\True}{\lnot\lexists[y][\formula{B}(y)]}, just=\TAss, checked
        [\sFmla{\False}{\lexists[y][\formula{B}(y)]},
          just={\TRule{\True}{\lnot}[3]}
          [\sFmla{\False}{\lforall[x][\formula{A}(x)]},
            just={\TRule{\True}{\lif}[2]},checked
            [\sFmla{\False}{\formula{A}(b)},
              just={\TRule{\False}{\lforall}[5]}
              [\sFmla{\True}{\formula{A}(b)},
                just={\TRule{\True}{\lforall}[1]},close
              ]
            ]
          ]
          [\sFmla{\True}{\lexists[y][\formula{B}(y)]},
            just={\TRule{\True}{\lif}[2]},checked
            [\sFmla{\True}{\formula{B}(c)},
              just={\TRule{\True}{\lexists}[5]}
              [\sFmla{\False}{\formula{B}(c)},
                just={\TRule{\False}{\lexists}[4]},close
              ]
            ]
          ]
        ]
      ]
    ]
  ]
\end{oltableau}
\end{ex}

\begin{prob}
కిందివాటికి సంవృత !!{tableau}s ఇవ్వండి:
\begin{enumerate}
\item $\sFmla{\False}{(\lforall[x][!A(x)] \land \lforall[y][!B(y)])
\lif \lforall[z][(!A(z) \land !B(z))]}$.
\item $\sFmla{\False}{(\lexists[x][!A(x)] \lor \lexists[y][!B(y)])
\lif \lexists[z][(!A(z) \lor !B(z))]}$.
\item $\sFmla{\True}{\lforall[x][(!A(x) \lif !B)]},
\sFmla{\False}{\lexists[y][!A(y)] \lif !B}$.
\item $\sFmla{\True}{\lforall[x][\lnot !A(x)]},
\sFmla{\False}{\lnot\lexists[x][!A(x)]}$.
\item $\sFmla{\False}{\lnot\lexists[x][!A(x)] \lif \lforall[x][\lnot
!A(x)]}$.
\item $\sFmla{\False}{\lnot\lexists[x][\lforall[y][((!A(x,y) \lif
\lnot !A(y,y)) \land (\lnot !A(y,y) \lif !A(x,y)))]]}$.
\end{enumerate}
\end{prob}

\begin{prob}
కిందివాటికి సంవృత !!{tableau}s ఇవ్వండి:
\begin{enumerate}
\item $\sFmla{\False}{\lnot\lforall[x][!A(x)] \lif \lexists[x][\lnot!A(x)]}$.
\item $\sFmla{\True}{(\lforall[x][!A(x)] \lif !B)}, \sFmla{\False}{\lexists[y][(!A(y) \lif !B)]}$.
\item $\sFmla{\False}{\lexists[x][(!A(x) \lif \lforall[y][!A(y)])]}$.
\end{enumerate}
\end{prob}

\end{document}
